\begin{pastexam}[label={px:nchu-113-dsa-26}]{中興資工所甲組 113 學年度 申論題 PART 2 第 2 題}
（問答）

(5\%) Assume P$\neq$NP. For each of the following problems, decide whether it is a P-problem, or an NP-hard
(or NP-complete) problem, or neither. Please answer 1 if it is a P-problem, answer 2 if it is an NP-hard (or
NP-complete) problem, and answer 3 if it is neither 1 nor 2.
\begin{enumerate}[label=(\arabic*)]
\item Graph Coloring: Given a graph, assign colors to each vertex so that no two adjacent vertices share the
same color, using a fixed number of colors.
\item Matrix Multiplication: Given two matrices, compute their product.
\item Traveling Salesman Problem: Given a list of cities and the distances between each pair, find the
shortest possible route visiting each city exactly once and returning to the origin city.
\item Determining if a Number is Prime: Given an integer, determine whether it is a prime number.
\item Sudoku: Given a partially filled grid of an arbitrary size, fill in the grid so that each row, column,
and square subgrid (of size determined by the grid size) contains all unique digits. Consider the decision
version of this problem, which asks whether a valid filling exists.
\end{enumerate}
\end{pastexam}

\begin{pastsoln}
答案：(1) 2　(2) 1　(3) 2　(4) 1　(5) 2。

在 $\mathrm{P} \ne \mathrm{NP}$ 的假設下，NP-hard 的問題不在 P 裡，所以 1 和 2 互斥（資結書定理~28.5）。
\begin{enumerate}[label=(\arabic*)]
\item \textbf{2}。 「用固定的 $k$ 種顏色著色」：$k \ge 3$ 時是 NP-complete（$3$-COLOR，資結書第~28.5~節的
表），題目的意思是一般的 $k$。 （$k = 2$ 是例外：就是判斷二分圖，用 BFS 在 $O(n + m)$ 時間解出，屬於 P；
$k$ 當作輸入的一部分時也是 NP-complete。）
\item \textbf{1}。 兩個 $n \times n$ 矩陣相乘，照定義三層迴圈 $\Theta(n^3)$（Strassen 更快，資結書第~17.3.4~節），
是多項式時間。 嚴格說它是「算出答案」而不是判定問題，但對應的判定版本（「乘積的第 $(i, j)$ 格是不是
$v$」）在多項式時間內可解，所以歸為 P。
\item \textbf{2}。 找最短的旅行推銷員迴路是最佳化問題；它的判定版本「有沒有總長 $\le L$ 的迴路」是
NP-complete（資結書定理~28.12），能解最佳化版本就能解判定版本，所以最佳化版本是 NP-hard。 它本身不是
判定問題，所以說 NP-hard 而不說 NP-complete。
\item \textbf{1}。 判斷質數在 P 裡：AKS 演算法（2002）的時間是輸入位數 $\lg N$ 的多項式。 注意輸入大小是
位數，從 $2$ 試除到 $\sqrt{N}$ 要 $\Theta(\sqrt{N}) = \Theta(2^{(\lg N)/2})$ 次，是指數時間，不能拿來說明它在
P 裡。 （質因數分解則不知道是否在 P 裡，但這題問的是判斷質數。）
\item \textbf{2}。 $n^2 \times n^2$ 的廣義數獨判定問題是 NP-complete：給一個填法，檢查每列、每行、每個小方格
都不重複只要多項式時間，所以在 NP；已知可以從 NP-complete 的問題（例如拉丁方陣補完）歸約過來（Yato
與 Seta，2003）。 一般的 $9 \times 9$ 數獨大小固定，是常數時間，題目特別說 arbitrary size 就是要排除這種情形。
\end{enumerate}
沒有一題的答案是 3。
\end{pastsoln}
